\newif\ifglobal

%%%%%%%%%%%%%%%%%%%%%%%
%%% Compile to view %%%
%%%%%%%%%%%%%%%%%%%%%%%

%%%%%%%%%%%%%%%%%%%%%%%%%%%%%%%%%%%%%%%%%%%%%%%%%%%
%%% Readme:                                     %%%
%%% https://www.overleaf.com/read/bwdjvfjksvjy/ %%%
%%%%%%%%%%%%%%%%%%%%%%%%%%%%%%%%%%%%%%%%%%%%%%%%%%%

% >>> M?: marks a module setting number or important notice

% >>> M4: Set {./relative-path} to external files for this module <<<
\def \modulefiles {../ESP32-P4/19-WDT}
% To add an external file, use:
% (Replace '---' with actual filename)
% '\input{\modulefiles/tables/---.tex}' for tables
% '\includegraphics{\modulefiles/figures/---}' for figures
% '\subfile{\modulefiles/---}' for register files

% >>> M1: This module is in global project? <<<
% 'YES' - uncomment; 'NO' - comment out
\globaltrue

\ifglobal
    \documentclass[main\_\_CN.tex]{subfiles}

\else
    % Font "FandolSong-Regular" used by ctex does not have GSUB table.
    % It causes warnings that do not affect compilation and can be ignored.
    \PassOptionsToPackage{quiet}{fontspec}

    \documentclass[openany, 10pt]{book}

    % >>> M2: In ./chip-spec-settings/preamble-trm-module, also find
    %         settings for visibility of todo notes, watermark <<<
    \usepackage{preamble-trm-module}


    % Enable Chinese typesetting
    \usepackage{ctex}

    % >>> M3: Temporary labels in Chinese
    %         you can add your own labels there <<<
    \myexternaldocument{temp-labels-cn}

\fi %\ifglobal


\setcounter{secnumdepth}{4}


% >>> M5: If this module needs any updates in
% './chip-spec-settings/preamble-trm-module.sty'
% (include additional package, etc.), contact your module PM
% or create the file './chip-spec-settings/changelog.tex' make the
% necessary updates yourself and document them in this file <<<

\raggedright


\begin{document}


% Sets language into which pre-defined bits of text are translated
\selectlanguage{chinese}

% Do not delete this block
\ifglobal
% Add the following front matter if
% this module is outside of global project
\else
    \InputIfFileExists{readme}{}{}
    {\let\clearpage\relax \listoftodos}
    {\let\clearpage\relax \listofdonetodos}
    \tableofcontents
    %\listoftables
    %\listoffigures
\fi


%%%%%%%%%%%%%%%%%%%%%%%%%
%%% TRM Module Begins %%%
%%%%%%%%%%%%%%%%%%%%%%%%%

\hypertarget{wdt}{}
\chapter{看门狗定时器 (WDT)}
\label{mod:wdt}

\section{概述}

看门狗定时器是一种硬件定时器，用于监视系统运行状态。软件必须定期喂狗（复位看门狗定时器）。若系统出现故障、死锁、程序卡住或其他异常情况导致软件未及时喂狗，看门狗定时器将超时并触发中断或复位，将系统导向至已知状态进行异常处理或重新启动。因此，看门狗定时器有助于检测、处理系统或软件的错误行为。

如图 \ref{fig:wdt-timers-overview} 所示，\chipname{} 中有三个数字看门狗定时器：章节 \ref{mod:timg} \textit{\nameref{mod:timg}} 描述的两个定时器组中各有一个主系统看门狗定时器，缩写为 MWDT，LP 系统中有一个 RTC 看门狗定时器，缩写为 RWDT。数字看门狗在运行期间会经历四个阶段（除非看门狗按时喂狗或者处于关闭状态），每个阶段均可配置单独的超时时间和超时动作，其中 MWDT 支持中断、HP CPU 复位和 HP 内核复位三种超时动作，RWDT 支持中断、HP CPU 复位、HP 内核复位和系统复位四种超时动作（详见章节 \ref{sec:wdt-timeout} \textit{\nameref{sec:wdt-timeout}}）。每个阶段的超时时间都可单独设置。

在 SPI Boot 模式下，RWDT 和定时器组 0 的 MWDT 会默认使能，以检测引导过程中发生的错误，并恢复运行。

\chipname{} 中还有一个模拟看门狗定时器——超级看门狗 (SWD)。超级看门狗是模拟域的超低功耗电路，可以防止系统在数字电路异常状态下运行，并在必要时复位系统（即超时动作中的系统复位）。

\begin{figure}[H]
    \centering
    \includegraphics[width=0.5\textwidth]{\modulefiles/figures/wdt-overview-lp}
    \caption{看门狗定时器概览}
    \label{fig:wdt-timers-overview}
\end{figure}

请注意，MWDT 寄存器部分详见章节 \ref{mod:timg} \textit{\nameref{mod:timg}}，RWDT 和 SWD 寄存器部分详见 \ref{sec:wdt-reg-summ} \textit{\nameref{sec:wdt-reg-summ}}。


\section{数字看门狗定时器}
\subsection{主要特性}
看门狗定时器具有如下特性：

\begin{itemize}
    \item 四个阶段，每个阶段都可配置超时时间和超时动作
    \item 超时动作
        \begin{itemize}
            \item MWDT：中断、HP CPU 复位、HP 内核复位
            \item RWDT：中断、HP CPU 复位、HP 内核复位、系统复位
        \end{itemize}
    \item 阶段 0 flash 启动保护（SPI Boot 模式）：
        \begin{itemize}
            \item MWDT0：超时触发 HP 内核复位
            \item RWDT：超时触发系统复位
        \end{itemize}
    \item 写保护，使能时寄存器仅可读取
    \item 32 位超时计数器
    \item 时钟源：
        \begin{itemize}
            \item MWDT：PLL\_F80M\_CLK、RC\_FAST\_CLK 或 XTAL\_CLK
            \item RWDT：LP\_DYN\_SLOW\_CLK
        \end{itemize}
\end{itemize}

\subsection{功能描述}

\begin{figure}[H]
    \centering
    \includegraphics[width=1.0 \textwidth]{\modulefiles/figures/761wdt.pdf}
    \caption{\chipname{} 的数字看门狗定时器}
    \label{fig:wdt-timers}
\end{figure}
图 \ref{fig:wdt-timers} 为 \chipname{} 数字系统中的三个看门狗定时器。

\subsubsection{时钟源与 32 位计数器}

看门狗定时器的核心是一个 32 位计数器。

以 MWDT0 为例：

\begin{itemize}
    \item MWDT0 可通过设置 \hyperref[regdesc:HPSYSTIMERGROUP0WDTCLKCONFREG]{HP\_SYS\_CLKRST\_PERI\_CLK\_CTRL20\_REG} 寄存器的 \hyperref[fielddesc:HPSYSTG0WDTCLKSEL]{HP\_SYS\_TIMERGRP0\_WDT\_SRC\_SEL} 字段选择 PLL\_F80M\_CLK 时钟、RC\_FAST\_CLK 或 XTAL\_CLK（外部时钟）作为时钟源。

    \item 将 \hyperref[regdesc:HPSYSTIMERGROUP0WDTCLKCONFREG]{HP\_SYS\_CLKRST\_PERI\_CLK\_CTRL20\_REG} 寄存器的 \hyperref[fielddesc:HPSYSTG0WDTCLKEN]{HP\_SYS\_TIMERGRP0\_WDT\_CLK\_EN} 字段置 1 开启时钟，清零关闭时钟。时钟经由可配置的 16 位预分频器分频。 详见章节 \ref{mod:resclk} \textit{\nameref{mod:resclk}} 的表 \ref{tab:resclk-cpu-clk-src}。

    \item MWDT0 的 16 位预分频器可通过 \hyperref[regdesc:TIMGWDTCONFIG1REG]{TIMG\_WDTCONFIG1\_REG} 寄存器的 \hyperref[fielddesc:TIMGWDTCLKPRESCALE]{TIMG\_WDT\_CLK\_PRESCALE} 字段配置。 \hyperref[fielddesc:TIMGWDTDIVCNTRST]{TIMG\_WDT\_DIVCNT\_RST} 字段置位时，预分频器复位，可立即重新配置。
\end{itemize}

RWDT 直接将 RTC 慢速时钟 LP\_DYN\_SLOW\_CLK（详见章节 \ref{mod:resclk} \textit{\nameref{mod:resclk}}）用作时钟源。

MWDT 和 RWDT 看门狗可分别通过设置 \hyperref[fielddesc:TIMGWDTEN]{TIMG\_WDT\_EN} 和 \hyperref[fielddesc:RTCWDTEN]{RTC\_WDT\_EN} 字段使能。看门狗使能后，其内部 32 位计数器的值会在每个时钟源周期内累加 1，直到达到该阶段的超时时间（即在该阶段发生超时）。如发生超时，计数器的值会重置为 0，同时看门狗进入下一阶段。如果软件成功喂狗，看门狗定时器会回到阶段 0，并将计数器的值重置为 0。软件向 \hyperref[regdesc:TIMGWDTFEEDREG]{TIMG\_WDTFEED\_REG} 寄存器内写入任意值，便可为 MDWT  喂狗；向 \hyperref[fielddesc:RTCWDTFEED]{RTC\_WDT\_FEED} 写 1，便可为 RWDT 喂狗。


\subsubsection{阶段与超时动作}\label{sec:wdt-timeout}

定时器在各阶段可以配置不同的超时时间和对应的超时动作。某一阶段超时会触发对应的超时动作，同时计数器的值被重置为 0，看门狗进入下一阶段。

MWDT 和 RWDT 有四个阶段（称为阶段 0 至阶段 3）。看门狗定时器会循环工作（即从阶段 0 至阶段 3，再回到阶段 0）。

MWDT 每个阶段的超时时间可用 \hyperref[regdesc:TIMGWDTCONFIG2REG]{TIMG\_WDTCONFIG\textcolor{red}{\it{i}}\_REG}（i 的范围是 2 到 5）寄存器配置，RWDT 的超时时间可用 \hyperref[fielddesc:RTCWDTSTG3HOLD]{RTC\_WDT\_STG\textcolor{red}{\it{j}}\_HOLD}（j 的范围是 0 到 3）字段配置。

值得注意的是，RWDT 在阶段 0 的超时时间 (T\textsubscript{hold0}) 受 eFuse 寄存器 \hyperref[regdesc:EFUSERDREPEATDATA0REG]{EFUSE\_RD\_REPEAT\_DATA0\_REG} 的 \hyperref[fielddesc:EFUSEWDTDELAYSEL]{EFUSE\_WDT\_DELAY\_SEL} 字段和 \hyperref[fielddesc:RTCWDTSTG0HOLD]{RTC\_WDT\_STG0\_HOLD} 字段共同影响，关系如下：

\begin{displaymath}
T_{hold0} = \hyperref[fielddesc:RTCWDTSTG0HOLD]{RTC\_WDT\_STG0\_HOLD}   << (\hyperref[fielddesc:EFUSEWDTDELAYSEL]{EFUSE\_WDT\_DELAY\_SEL} + 1)
\end{displaymath}

其中，$<<$ 为左移运算符。例如，\hyperref[fielddesc:RTCWDTSTG0HOLD]{RTC\_WDT\_STG0\_HOLD} 配置为 100，\hyperref[fielddesc:EFUSEWDTDELAYSEL]{EFUSE\_WDT\_DELAY\_SEL} 为 1，则 T\textsubscript{hold0} 为 400 个时钟周期。

如某个阶段超时，下列超时动作之一将会执行：

\begin{table}[H]
\centering
\caption{超时动作}
\label{tab:wdt-timeout-action}

\begin{tabular}{|l|m{12cm}|}
\hline
\rowcolor{lightgray}
    \textbf{超时动作} & \textbf{描述} \\\hline


    中断 &  触发中断 \\ \hline
    HP CPU 复位 &  复位 HP CPU0 和 HP CPU1\\ \hline
    HP 内核复位 & HP 内核复位包括 HP CPU0、HP CPU1、HP 外设、HP GPIO 等\\ \hline
    系统复位 &  复位包括低功耗系统在内的整个数字系统 \\ \hline
    关闭 &  对系统不产生影响\\ \hline

\end{tabular}
\end{table}

MWDT 所有阶段的超时动作均在 \hyperref[regdesc:TIMGWDTCONFIG0REG]{TIMG\_WDTCONFIG0\_REG} 寄存器中配置。RWDT 的超时动作可在\\ \hyperref[regdesc:RTCWDTCONFIG0REG]{RTC\_WDT\_CONFIG0\_REG} 寄存器配置。

\subsubsection{写保护}
看门狗定时器对于检测和处理系统或软件错误而言至关重要，不应轻易关闭（例如，因写寄存器位置错误而误将看门狗关闭）。因此，MWDT 和 RWDT 引入写保护机制，防止看门狗因无意的写操作而被关闭或篡改。

写保护机制通过每个看门狗定时器的写密钥字段运行（MWDT 看门狗使用 \hyperref[fielddesc:TIMGWDTWKEY]{TIMG\_WDT\_WKEY}，RWDT 看门狗使用 \hyperref[fielddesc:RTCWDTWKEY]{RTC\_WDT\_WKEY}）。必须向看门狗定时器的写密钥字段写入 0x50D83AA1，才能修改其它看门狗寄存器。如果写密钥字段的值不是 0x50D83AA1，任何试图向看门狗定时器寄存器（除了向写密钥字段本身）写值的操作都会被忽略。推荐按以下步骤访问看门狗定时器：

\begin{enumerate}
\item 将 0x50D83AA1 写入看门狗定时器的写密钥字段，关闭写保护。
\item 根据需要修改看门狗，如喂狗或改变配置。
\item 向看门狗定时器的写密钥字段上写入除 0x50D83AA1 以外的任意值，重新使能写保护。
\end{enumerate}

\subsubsection{flash 引导保护}
在 SPI boot 模式下，MWDT0 和 RWDT 会默认使能。MWDT0 的阶段 0 的默认超时动作为 HP 内核复位（复位主系统）。RWDT 的阶段 0 超时动作为系统复位（复位主系统和 RTC）。完成 ROM 阶段的 SPI boot 引导后，应将 \hyperref[fielddesc:TIMGWDTFLASHBOOTMODEN]{TIMG\_WDT\_FLASHBOOT\_MOD\_EN} 和 \hyperref[fielddesc:RTCWDTFLASHBOOTMODEN]{RTC\_WDT\_FLASHBOOT\_MOD\_EN} 位清零，分别关闭 MWDT0 和 RWDT 的 flash 引导保护。然后，软件可以配置 MWDT0 和 RWDT。

\section{超级看门狗定时器}\label{sec:wdt-swdt}
超级看门狗 (SWD) 是模拟域的超低功耗电路，可以防止系统在数字电路异常状态下运行，并在必要时复位系统（即超时动作中的系统复位）。SWD 包含一个看门狗电路，需在每个超时阶段（约不足一秒）至少喂狗一次。该电路会在看门狗超时时间约 100 ms 之前发送 WD\_INTR 信号提醒系统喂狗。

如果系统不回应 SWD 的喂狗请求，看门狗超时，SWD 会产生系统电平信号 SWD\_RSTB，复位芯片上的整个数字电路（即超时动作中的系统复位）。

SWD 的时钟源固定，不可选择。
\subsection{主要特性}
SWD 具有如下特性：
\begin{itemize}
    \item 超低功耗
    \item 用中断提醒 SWD 即将超时
    \item 软件有多种专用的方法喂 SWD，让 SWD 监控整个操作系统的工作状态
\end{itemize}

\subsection{SWD 控制器}
\subsubsection{结构}

\begin{figure}[H]
    \centering
    \includegraphics[width=0.9\textwidth]{\modulefiles/figures/swd_pic_thicker}
    \caption{SWD 控制器结构}
    \label{fig:wdt}
\end{figure}
\subsubsection{工作流程}
正常状态下：
\begin{itemize}
    \item SWD 控制器收到 SWD 的喂狗请求。
    \item SWD 控制器向主 CPU 发送中断。
    \item 主 CPU 通过置位 \hyperref[fielddesc:RTCWDTSWDFEED]{RTC\_WDT\_SWD\_FEED} 直接喂狗。
    \item CPU 喂狗时，需要先向 \hyperref[fielddesc:RTCWDTSWDWKEY]{RTC\_WDT\_SWD\_WKEY} 写 0x50D83AA1 关闭 SWD 控制器的写保护。这样做可以防止系统在数字电路异常状态下运行时误喂 SWD。
    \item 如将 \hyperref[fielddesc:RTCWDTSWDAUTOFEEDEN]{RTC\_WDT\_SWD\_AUTO\_FEED\_EN} 置 1，SWD 控制器也可配置为在不需要 CPU 干预的情况下喂 SWD。
\end{itemize}

复位后：
\begin{itemize}
    \item 可查看 \hyperref[fielddesc:LPAONCLKRSTLPCORERESETCAUSE]{LP\_AONCLKRST\_LPCORE\_RESET\_CAUSE}[4:0] 获知 HP CPU 复位原因。\\
	如 \hyperref[fielddesc:LPAONCLKRSTLPCORERESETCAUSE]{LP\_AONCLKRST\_LPCORE\_RESET\_CAUSE}[4:0] == 0x12，则表示上一次复位的原因是 SWD 复位。
	\item 置位 \hyperref[fielddesc:RTCWDTSWDRSTFLAGCLR]{RTC\_WDT\_SWD\_RST\_FLAG\_CLR} 清除 SWD 复位标志。
\end{itemize}



%%%%%%%%%%%%%%%%%%%%%
%%% Interrupts %%%
%%%%%%%%%%%%%%%%%%%%%

\section{中断}

看门狗定时器中断，请前往章节 \ref{mod:timg} \textit{\nameref{mod:timg}} 的第 \ref{sec:timg-interrupt} 节 \textit{\nameref{sec:timg-interrupt}} 查看。

% >>> If your module has no registers, delete sample text
%       for Sections Register Summary and Registers below <<<

%%%%%%%%%%%%%%%%%%%%%%%%
%%% Register Summary %%%
%%%%%%%%%%%%%%%%%%%%%%%%

% >>> Update the sample text below and insert
%     the info on register summary into the subfile <<<

\newpage
\hypertarget{wdt-reg-summ}{}
\section{寄存器列表}
\label{sec:wdt-reg-summ}

MWDT 寄存器是定时器组模块的一部分，在章节 \ref{mod:timg} \textit{\nameref{mod:timg}} 的第 \ref{sec:timg-reg-summ} 节 \textit{\nameref{sec:timg-reg-summ}} 中有详细描述。


本小节包含 RWDT 和 SWD 相关寄存器介绍，其地址均为相对于 RTC\_WDT 基地址的地址偏移量（相对地址），具体基地址请见章节 \ref{mod:sysmem} \textit{\nameref{mod:sysmem}} 中的表 \ref{tab:sysmem-base-address}。

请查看章节 \hyperref[glossary-access-types]{\textit{寄存器的访问类型}}，了解“访问”列缩写的含义。


\subfile{\modulefiles/19-WDT-rwdt-reg-summ__CN}


%%%%%%%%%%%%%%%%%
%%% Registers %%%
%%%%%%%%%%%%%%%%%

% >>> Update the sample text below and insert
%      the info on registers into the subfile <<<

\newpage
\section{寄存器}

MWDT 寄存器是定时器组模块的一部分，在章节 \ref{mod:timg} \textit{\nameref{mod:timg}} 的第 \ref{sec:timg-reg-summ} 节 \textit{\nameref{sec:timg-reg-summ}} 中有详细描述。


本小节包含 RWDT 和 SWD 相关寄存器介绍，其地址均为相对于 RTC\_WDT 基地址的地址偏移量（相对地址），具体基地址请见章节 \ref{mod:sysmem} \textit{\nameref{mod:sysmem}} 中的表 \ref{tab:sysmem-base-address}。


\subfile{\modulefiles/19-WDT-rwdt-reg__CN}

\end{document}

